\section{Finite comparisons and the structural height coefficient}
\label{st:section}

The output of this section is a coefficient $a(k)$ determined by
Gaussian limits on free rectangles.  Finite lattice comparisons also
give the gap $a(k)=0$ or $a(k)\ge8\pi$, and localization on the zero
branch.  Section~\ref{an:section} will relate this coefficient to the
physical roughness threshold and to the field on other domains.

Throughout this section $J$ has the hypotheses in the introduction, $v=v_J$,
\begin{equation}\label{st:normalization}
 A_0=-\Delta_J/v^2,\qquad k=\beta/v^2,\qquad h=\sigma\in2\pi\Z.
\end{equation}
An unoriented interaction edge has conductance $c=(|J|v^2)^{-1}$; parallel
edges are counted with their multiplicities. Thus
$(g,A_0g)=\sum_e c(\nabla_e g)^2$, and the Fourier symbol of $A_0$ is
$|p|^2+O_J(|p|^4)$. Write $R_J=\max_{z\in J}|z|_\infty$.
For a finite rectangle $B$ retain only edges whose two endpoints belong to
$B$, and denote the resulting Laplacian and its inverse on constants'
orthogonal complement by $A_B$ and $G_B=A_B^+$. A free height law is the
probability law modulo common translations by $2\pi\Z$, with weight
$\exp\{-(h,A_Bh)/(2k)\}$. Pinning one vertex is one choice of representative.
All pairings on a finite graph are counting-measure pairings.

Here is a useful convention concerning errors. An edge representation
$X=\sum_e w_e\nabla_e h$ has \emph{flow cost}
$\sum_e w_e^2/c$. For a zero-sum vertex profile $f$,
\begin{equation}\label{st:flow}
 (f,G_Bf)=\min_{\operatorname{div}w=f}\sum_e w_e^2/c.
\end{equation}
Indeed the minimizer is $w=c\nabla G_Bf$; its difference from every other
admissible flow is orthogonal to gradients in the scalar product
$\sum_e w_e w'_e/c$. In particular a test of flow cost $\delta$ has variance
at most $k\delta$ and all its exponential moments have the corresponding
Gaussian bound, by Lemma~\ref{st:finite-gaussian}. Nearest-neighbor edges,
which are present by hypothesis, may always be used to construct trial flows.

\subsection{Finite lattice Gaussian inequalities}

We start from the lattice inequalities of
\citet[Theorem~2.1, Corollary~3.2, and Propositions~3.3 and~4.2]{RSD2017}
and derive the comparison and limiting forms used below, including
proofs to specify the directions.
For applications of these inequalities to integer-height depinning
and graph surgery, see \citet{AizenmanHarelPeledShapiro2022}.

\begin{lemma}[Finite Gaussian comparisons]\label{st:finite-gaussian}
Let $\mathcal L$ be a lattice in a finite-dimensional Euclidean space $E$,
let $V=\operatorname{span}\mathcal L$, and let $P_V$ be orthogonal projection
onto $V$. Suppose $T\ge0$ is symmetric and $T_V=P_VT|_V$ is positive
definite. For $x\in E$ set
$w_T(x)=\sum_{u\in\mathcal L+x}\exp\{-(u,Tu)/2\}$.
Then
\begin{equation}\label{st:lattice-inequality}
 w_T(0)^2w_T(x+y)w_T(x-y)\ge w_T(x)^2w_T(y)^2.
\end{equation}
The covariance of the Gaussian on $\mathcal L+x$ is at least that on
$\mathcal L$ at the same precision, in positive-semidefinite order; the
same lower bound holds after a real linear tilt of the shifted law. For
each fixed ambient displacement $x$, the ratio $w_{T+tB}(x)/w_{T+tB}(0)$
is nonincreasing in $t\ge0$ whenever $B\ge0$. For the centered Gaussian
$H$ on $\mathcal L$, every real linear test $X=(\ell,H)$ satisfies
\begin{equation}\label{st:mgf-sandwich}
 \tfrac12\Var(X)\le\log\E e^X
 \le\tfrac12(\ell_V,T_V^{-1}\ell_V),\qquad \ell_V=P_V\ell.
\end{equation}
Centered variances and real moment generating functions decrease when
precision increases. For any two real linear tests,
\begin{equation}\label{st:paired-cumulant}
 \kappa(X,X,Y,Y)
 =\E X^2Y^2-\E X^2\E Y^2-2(\E XY)^2\ge0.
\end{equation}
The centered characteristic function is strictly positive, and its logarithm
has Hessian everywhere at least its Hessian at zero.
These centered conclusions apply after added centered pins or equalities.
For free graphs they apply to real tests annihilating every released
constant. The centered characteristic function is strictly positive for
characters invariant under every released translation, and its logarithmic
Hessian comparison holds along directions annihilating those constants.
The probability of a specified collection of centered lattice equalities
increases under a positive-semidefinite increase of precision.
\end{lemma}

\begin{proof}
For $d\in\mathcal L/2\mathcal L$ put
$z_d(s)=\sum_{u\in d+2\mathcal L}e^{-(u+s,T(u+s))/4}$.
Coercivity on $V$ makes these sums and their derivatives locally uniformly
convergent for ambient shifts.
The bijection $(u,v)\mapsto(u+v,u-v)$ gives
\[
 w_T(x)w_T(y)=\sum_d z_d(x+y)z_d(x-y),\qquad
 \sum_d z_d(s)^2=w_T(0)w_T(s).
\]
Cauchy--Schwarz proves Equation~\eqref{st:lattice-inequality}.
Taking logarithms and differentiating twice in $y$ at zero gives
$D^2\log w_T(x)\ge D^2\log w_T(0)$. Direct differentiation of the sum gives
$D^2\log w_T(x)=-T+T\Cov_x(H)T$. Both covariances are supported on $V$.
Restricting the Hessian inequality to directions in $V$ and inverting $T_V$
there proves covariance minimality.

For $x_T=T_V^{-1}P_VTx$ and $u\in V$, completion of the square gives
\[
 (u+x,T(u+x))=(u+x_T,T_V(u+x_T))+c_T(x),\qquad
 c_T(x)=(x,Tx)-(x_T,T_Vx_T)\ge0.
\]
For fixed $x$ and $T_t=T+tB$, $B\ge0$, differentiation in the ambient
space gives
\[
 \frac{d}{dt}\log\frac{w_{T_t}(x)}{w_{T_t}(0)}
 =-\tfrac12\tr\bigl[B(\Cov_x(H)-\Cov_0(H)
                         +\E_xH\otimes\E_xH)\bigr]\le0.
\]
This derivative keeps the ambient displacement fixed and includes the
$T_t$-dependence of both $x_{T_t}$ and the normal scalar $c_{T_t}(x)$.
A real tilt by $\ell$ changes the effective displacement on $V$ to
$x_T-T_V^{-1}\ell_V$. Its covariance therefore still dominates the centered
covariance. For the centered law, integrating the Hessian of its logarithmic
moment generating function on the segment from $0$ to $\ell$ proves the
first inequality in Equation~\eqref{st:mgf-sandwich}; symmetry makes the
first derivative at zero vanish. Differentiating this logarithmic moment
generating function with respect to $T_t$ gives the same nonpositive
expression with the tilted covariance and mean. Differentiating the resulting
centered comparison twice at zero proves the variance comparison. The second
derivative at $s=0$ of $\Var_{sY}(X)\ge\Var_0(X)$ proves
Equation~\eqref{st:paired-cumulant}.

Let $\mathcal L^*$ be the dual lattice in $V$. Poisson summation on $V$ gives
\[
 w_T(x)=e^{-c_T(x)/2}C_{T_V}\sum_{q\in\mathcal L^*}
 e^{-2\pi^2(q,T_V^{-1}q)}e^{2\pi i(q,x_T)}\le w_T(0).
\]
Completion of the square now proves the upper bound in
Equation~\eqref{st:mgf-sandwich}. For the centered characteristic function,
the same summation gives the strictly positive formula
\[
 \E e^{i(\ell,H)}
 =\frac{\displaystyle\sum_{q\in\mathcal L^*}
 e^{-\frac12(2\pi q-\ell_V,T_V^{-1}(2\pi q-\ell_V))}}
 {\displaystyle\sum_{q\in\mathcal L^*}e^{-2\pi^2(q,T_V^{-1}q)}}>0.
\]
It is a shifted-to-centered partition ratio on the dual lattice.
Equation~\eqref{st:lattice-inequality} there proves the assertion about its
logarithmic Hessian.

A collection of centered equalities selects
$\mathcal L_E=\mathcal L\cap\bigcap_j\ker\ell_j$, a lattice in its span.
The restricted precision is coercive, so the preceding argument applies
directly to its conditional centered law. In particular the displayed dual
sum proves strict characteristic positivity there. This law is also the
dominated limit of the additional precisions
$t\sum_j\ell_j\otimes\ell_j$, $t\uparrow\infty$, which transfers the centered
comparisons with the original law.

For a free graph choose one representative $h_{x_0}=0$ on each connected
component. Its representative lattice has a span on which the massless graph
precision is coercive. The same dual formula directly proves strict
characteristic positivity on that centered law, and hence for its
translation-invariant characters. To compare a graph with a cut graph,
first add $\eta\Id$.
On a connected component choose representatives $h_{x_0}=0$ and sum its
common translation $2\pi j$. Completion of the square shows that this sum,
divided by $(2\pi\eta|B|)^{-1/2}$, tends to a constant independent of the
representative and is uniformly bounded for $0<\eta\le\eta_0$.
The remaining representative weight is bounded by a constant times its
massless weight and converges pointwise to it. Dominated convergence,
also after any fixed real tilt annihilating constants, transfers the
comparison inequalities.
The same argument applies to bounded observations invariant under common
$2\pi$ translations, a fact used below for phases.
Finally, for an equality event $E$, differentiate its probability along
$T+tB$ before taking any infinite-penalty limit. Its logarithmic derivative
is $\frac12\tr B(\Cov(H)-\Cov(H\mid E))\ge0$ by the just-proved centered
pin comparison; both means vanish by inversion symmetry. This proves the
last assertion, and dominated limits give its free-graph versions whenever
the event is invariant under the released translations.
\end{proof}

For later use, a consequence of the last Hessian assertion is worth recording.
Let $\phi$ be the centered characteristic function in
Lemma~\ref{st:finite-gaussian}, let $s$ be a base profile, and let $Y$ have
bare variance at most $V$. On a free graph take $s$ with integer total on
each released component and $Y$ with zero total on each. Then
\begin{equation}\label{st:coefficient-transfer}
 \sqrt{-\log\phi(s+tY)}
 \le\sqrt{-\log\phi(s)}+|t|\sqrt{V/2}.
\end{equation}
To prove this, set $f(t)=\log\phi(s+tY)\le0$. The inequality $f''\ge-V$
and the upper bound $f\le0$, applied at the maximum of
$f(t)+f'(t)u-Vu^2/2$, imply $|f'|^2\le-2Vf$.
Integration, or regularization at zeros of $-f$, gives
Equation~\eqref{st:coefficient-transfer}. In particular, a positive lower
bound transfers between profiles at bounded bare distance in either direction.

\subsection{Saturation on free rectangles}

Let $B_m=\{0,\ldots,m-1\}^2$, $g_m(x)=x_1/m$, and
\begin{equation}\label{st:am}
 a_m(k)=\Var_{B_m,k}(h,A_{B_m}g_m).
\end{equation}
The energy $(g_m,A_{B_m}g_m)$ tends to one, by square symmetry and the
definition of $v$. Thus $0\le a_m\le k+o(1)$.

\begin{theorem}[All free-rectangle limits]\label{st:free-rectangle-limit}
There is a number $a=a(k)\in[0,k]$ such that $a_m(k)\to a$.
Let $B_n$ be discrete axis rectangles at mesh $n^{-1}$ whose endpoints
converge to those of a fixed nondegenerate rectangle $B$. For every finite
collection of smooth shifts $g$ on a neighborhood of $\overline B$,
\begin{equation}\label{st:rectangle-laplace}
 \log\E_{B_n,k}\exp\{(h,A_{B_n}g(\cdot/n))\}
 \longrightarrow \frac a2\int_B|\nabla g|^2.
\end{equation}
The corresponding tests converge jointly to centered Gaussian variables
with covariance $a\int_B\nabla g\cdot\nabla g'$; all joint moments converge.
For zero-total profiles $f_n$ whose scaled piecewise constant densities
$n^2f_n$ converge in $L^2(B)$ to a bounded piecewise continuous $f$ of
integral zero, the corresponding assertion is
\begin{equation}\label{st:rectangle-smear}
 \log\E_{B_n,k}e^{(h,f_n)}
 \longrightarrow \frac a2(f,(-\Delta_{B,N})^{-1}f)_{L^2(B)}.
\end{equation}
Here the inverse Neumann Laplacian is on integral-zero functions. These
statements hold along every such sequence of rectangles and jointly for
mixed finite lists of the displayed tests.
\end{theorem}

\begin{proof}
We give separately the cutting, folding, and distributional steps.

\emph{Cutting and the existence of $a$.}
Cut a square of side $M$ into $m$-squares and strips of width less than $m$.
In its affine observation $x_1/M$, the squared cost of omitted seam and
strip flows is at most $C_J(m^{-1}+m/M)$: there are at most
$C_J(M^2/m+mM)$ affected edges and each coefficient is at most $C_J/M$.
On the remaining squares the observation is $m/M$ times the cell
observation. Cutting increases variance, and the cut components are
independent. The triangle inequality in $L^2$ consequently gives
\begin{equation}\label{st:cut-am}
 \sqrt{a_M}\le\sqrt{a_m}+C_J\sqrt{k}(m^{-1}+m/M)^{1/2}.
\end{equation}
First let $M\to\infty$, then take $m\to\infty$ along a liminf subsequence.
This proves convergence. Notice also that each $a_m(k)$ is continuous
for $k>0$, by locally uniform convergence of its defining Gaussian sums,
and Equation~\eqref{st:cut-am} is uniform when $k$ ranges over a compact
subset of $(0,\infty)$. In particular
$\sqrt{a(k)}\le\sqrt{a_m(k)}+C_J\sqrt{k/m}$.
Square symmetries make the covariance matrix of
the two affine tests scalar. Cutting a rectangle into small squares proves
the affine upper bound $a\int_B|\nabla g|^2$. For its reverse, let
$m_1(n),m_2(n)$ be the rectangle's side counts and set
$V_n(p)=\Var_{B_n}(h,A_{B_n}(p\cdot x/n))$. Tile a square of side $M$,
divisible by both side counts, with copies of this rectangle. The
normalized affine test $p\cdot x/M$ restricts on each cell to $n/M$
times the cell test, up to an irrelevant constant. There are
$M^2/(m_1(n)m_2(n))$ cells. Cutting and the seam-flow estimate, followed
by $M\to\infty$ at this fixed $n$, give
\[
 \sqrt a\,|p|
 \le\frac{n}{\sqrt{m_1(n)m_2(n)}}\sqrt{V_n(p)}
       +\frac{C_{J,k,p}}{\sqrt{\min(m_1(n),m_2(n))}}.
\]
Since $m_1(n)m_2(n)/n^2\to|B|$, this proves
$\liminf_nV_n(p)\ge a|p|^2|B|$ along every admissible rectangle sequence.
Together with the cutting upper bound it gives
\[
 \lim_n\Var_{B_n}(h,A_{B_n}g)=a\int_B|\nabla g|^2
\]
for every affine $g$ and every such sequence.

Cutting into cells of small fixed macroscopic diameter and replacing
$\nabla g$ by constants proves, for smooth or continuous piecewise smooth
$g$ with bounded piece derivatives and polygonal interfaces,
\begin{equation}\label{st:upper-energy}
 \limsup_n\Var(h,A_{B_n}g)\le a\int_B|\nabla g|^2.
\end{equation}
At a fixed partition the omitted edge count is $O(n)$ and coefficients
are $O(n^{-1})$, so their cost tends to zero. In cell interiors the error
cost is bounded by the $L^2$ modulus of continuity of the gradient, which
tends to zero as the partition is refined. The identical proof applies to
partial gradient observations restricted to polygonal pieces.

Call a shift $g$ saturated on $B$ if, for every admissible sequence
of discrete rectangles $B_n$ converging to $B$,
\[
 \lim_{n\to\infty}\Var_{B_n}(h,A_{B_n}g(\cdot/n))
       =a\int_B|\nabla g|^2.
\]
For a partial gradient observation use the same full-limit definition
with the energy integral over its retained region.

This class is linear. Fix a finite list of saturated shifts and any
admissible rectangle sequence. From any subsequence one can extract a
further subsequence on which their covariance matrix converges, by the
bare bounds. If $C$ is that limiting matrix and
$E_{ij}=a\int_B\nabla g_i\cdot\nabla g_j$, then
Equation~\eqref{st:upper-energy}, applied to every real linear combination,
gives $E-C\ge0$. Full convergence for each $g_i$ gives
$(E-C)_{ii}=0$ on this same subsequence. A positive-semidefinite matrix
with zero diagonal has all its entries zero. Thus $C=E$.
Since this applies to every initial subsequence, the covariance matrices
converge along the original sequence. In particular every finite linear
combination is saturated.

Saturation on a larger rectangle passes to an enclosed rectangle.
Partition the complement into finitely many rectangles, indexed together
with the inner rectangle by $i=0,\ldots,r$, and write
$E_i=\int_{B_i}|\nabla g|^2$. Vanishing seam-flow errors and cutting give
\[
 V_{\mathrm{outer},n}\le\sum_{i=0}^rV_{i,n}+o(1).
\]
The outer variance converges to $a\sum_iE_i$, while
$\limsup_nV_{i,n}\le aE_i$ for every piece. Consequently
\[
 \liminf_nV_{0,n}
 \ge a\sum_iE_i-\sum_{i=1}^r\limsup_nV_{i,n}
 \ge aE_0.
\]
Together with the upper bound this is full convergence on the inner
rectangle, on every admissible sequence. Strict enlargements suffice;
compatible rounding of the partition changes only vanishing-cost flows.

The class is also closed under gradient-energy approximation. If
$g_j$ are saturated and $\|\nabla(g_j-g)\|_{L^2(B)}\to0$, with convergence
of the corresponding discrete difference energies, then the bare bound
and the triangle inequality for standard deviations give
\[
 \limsup_n\left|
 \sqrt{\Var(h,A_{B_n}g)}-\sqrt a\,\|\nabla g\|_{L^2(B)}
 \right|
 \le(\sqrt k+\sqrt a)\|\nabla(g-g_j)\|_{L^2(B)}.
\]
Letting $j\to\infty$ proves full convergence for $g$. The same argument
shows that an observation changed by a flow of cost tending to zero
retains its full variance limit. These statements apply to the fixed
smooth and continuous piecewise smooth shifts used below, whose discrete
energies converge by Riemann summation.

\emph{The fold.}
We justify the FKG correlation input \citep{FKG1971} directly.
A quadratic attractive height
law has the lattice condition
$p(u\wedge v)p(u\vee v)\ge p(u)p(v)$, since it holds edge by edge for
$e^{-c(h_x-h_y)^2/(2k)}$. On a finite height box this condition implies
association by induction in the number of vertices. Conditional laws in
one fewer coordinate satisfy the same condition. Between two values of
the last coordinate the likelihood ratio of their conditional laws is
increasing. The induction hypothesis applied to that ratio proves that
the conditional expectation of every increasing function increases.
The conditional covariance decomposition and the one-variable identity
$\Cov(F(Z),G(Z))=\frac12\E[(F(Z)-F(Z'))(G(Z)-G(Z'))]\ge0$
complete the induction. Letting the box grow proves association in our
proper finite Gaussian law. This argument also proves association after
conditioning any set of vertices.

Suppose a rectangle and its discrete graph are invariant under reflection
across a site axis or diagonal. Pin a site of the reflection line and
condition its entire height trace. The conditional law remains invariant
under reflection and associated, including interactions crossing the line.
For a saturated shift odd about that line up to a constant, let $X_+$ and
$X_-$ be its observations on edges internal to the two open sides. Then
$X_-$ is the reflected negative of $X_+$ and their sum differs from the
full observation by a vanishing-cost flow. If $X_+$ can be changed by such
a flow to have nonnegative off-line vertex coefficients, conditional
association gives a nonpositive conditional covariance with $X_-$.
Their conditional means are negatives, so the unconditional covariance
is also nonpositive. Write $V_{\pm,n}=\Var(X_\pm)$ and
$E_\pm=\int_{B\cap\{\pm\text{ side}\}}|\nabla g|^2$.
The vanishing-cost modifications give
$\Cov(X_+,X_-)\le o(1)$, and the full observation differs from
$X_++X_-$ by a vanishing-cost flow. Its full variance convergence
therefore gives
\[
 a(E_++E_-)
 =\lim_n V_{\mathrm{full},n}
 \le\liminf_n(V_{+,n}+V_{-,n}).
\]
More directly, $V_{+,n}\ge V_{\mathrm{full},n}-V_{-,n}-o(1)$,
so the separate upper bound gives
\[
 \liminf_nV_{+,n}
 \ge a(E_++E_-)-\limsup_nV_{-,n}\ge aE_+.
\]
The same argument applies to the minus side. With
$\limsup_nV_{\pm,n}\le aE_\pm$, both partial variances converge
along the full sequence. The folded-shift observation differs from its
partial observation by a vanishing-cost flow, so it too has the required
full variance limit.

Here are the required coefficient checks and error bounds. For an axis
fold of $g=x_1$, pair each step $(z_1,z_2)$ with $(-z_1,z_2)$.
The coefficients of $A_{B_+}g$ vanish away from the two normal-face strips,
are nonpositive in the strip adjacent to the cut, and nonnegative at the
outer normal face. Tangential boundaries preserve this cancellation.
Move every cut-strip coefficient to the conditioned line along at most
$R_J+1$ nearest-neighbor edges. There are $O(n)$ coefficients of size
$O(n^{-1})$, and the paths have bounded overlap, giving cost $O(n^{-1})$.
Thus $(x_1-t)_+$ is saturated after coherent rounding of the reflection
line and restriction from an enclosing symmetric rectangle. Finite linear
combinations and gradient approximation give every smooth function of
either coordinate.

For the diagonal fold put $Y=y+t$ and use an enclosing square
$p\le x,Y\le q$. Take $g=F(x)-F(Y)$ with $F'''<0$ and
$F'(p),F'(q)>0$ on a slightly larger interval. In the interior of $x>Y$,
$A_Bg\ge0$: the symmetric second differences of $F$ decrease with their
center, and the step lists in the two coordinates coincide. At the outer
face $x=q$ the missing-step contribution has leading term a positive
constant times $F'(q)/n$; at $Y=p$ it has leading term a positive constant
times $F'(p)/n$. The tangential first-order terms cancel in pairs.
Uniform strict endpoint positivity absorbs the $O(n^{-2})$ remainder.
More precisely, in integer coordinates $P\le i,j\le Q$, the exceptional
sets can be taken as $D=\{0<i-j\le2R_J\}$ and
$K=\{Q-i\le R_J,\ j-P\le R_J\}$. Outside $D$ neither the lower $i$-face
nor the upper $j$-face truncates a step; outside $K$ the other two faces
cannot both truncate. This proves the stated sign everywhere outside
$D\cup K$. Move all coefficients in $D\cup K$ horizontally from $(i,j)$
to $(j,j)$ on the line. There are $O(n)$ short paths of length at most
$2R_J$ and $O(1)$ paths of length $O(n)$, with bounded overlap.
Their coefficients have size $O(n^{-1})$. If the overlap is $M$, their
cost is at most $(M/c)\sum_v d(v)^2|P_v|=O(n^{-1})$.
Define the opposite-side modification by reflected negative, so the
line terms cancel exactly. The fold criterion proves
saturation of
\begin{equation}\label{st:diagonal-fold}
 (F(x)-F(y+t))\ind_{\{x>y+t\}}.
\end{equation}
Every smooth $F$ on the fixed interval is a difference of two functions
with the strict properties just used: add a sufficiently large fixed
cubic with negative third derivative, and then a sufficiently large
positive linear function. Thus Equation~\eqref{st:diagonal-fold} is
saturated for arbitrary smooth $F$.

For clarity, coherent rounding preserves these full limits. Choose
$t_n=r_n/n$ with $r_n=\operatorname{round}(nt)$, and enclose the target
rectangle in the exactly symmetric lattice square
$P_n\le i\le Q_n$, $P_n-r_n\le j\le Q_n-r_n$.
Its reflection is $(i,j)\mapsto(j+r_n,i-r_n)$, and its scaled endpoints
can be chosen to converge. The odd seed $F(x)-F(y+t_n)$ differs from the
fixed saturated seed $F(x)-F(y+t)$ by discrete gradient energy
$O(|t_n-t|^2)$. For the folded functions $f_{t_n},f_t$ of
Equation~\eqref{st:diagonal-fold}, the squared difference energy is at most
$C(|t_n-t|+n^{-1}+|t_n-t|^2)$. Indeed the strip between their diagonals,
enlarged by $R_J/n$, meets $O(n^2|t_n-t|+n)$ edges, each carrying an
$O(n^{-1})$ difference increment; away from it, smooth derivative
differences are $O(|t_n-t|)$. Axis hinges satisfy the same estimate.
Thus the fold and restriction limits transfer to each fixed target
shift along every admissible rectangle sequence by the vanishing-cost
perturbation bound.

To see that these folds suffice, extend a smooth target $g$ slightly and
choose a smooth family $F_t$ of compact support in $t$, on all relevant
parameters, satisfying $F_t'(x)=-g_{xy}(x,x-t)$. Differentiating the integral
of Equation~\eqref{st:diagonal-fold} gives
\[
 \partial_x\partial_y\int
 (F_t(x)-F_t(y+t))\ind_{\{x>y+t\}}\,dt=-F_{x-y}'(x)=g_{xy}(x,y).
\]
The residual is a sum of two univariate functions. Riemann sums of this
integral converge in gradient $L^2$: translations of the bounded piecewise
smooth gradients are continuous in $L^2$, including across the moving
diagonal. The bare bound then transfers saturation to $g$.

\emph{Laplace transforms and smears.}
The lower bound in Equation~\eqref{st:rectangle-laplace} follows from
saturation and Equation~\eqref{st:mgf-sandwich}. For the upper bound, cut
into cells of side comparable to a small fixed $r$. For fixed $r$, the
discarded-flow observation $E_n$ has bare variance tending to zero.
H\"older with exponents $p>1$ and $p/(p-1)$ removes $E_n$ at a cost tending
to zero. Precision comparison bounds the remaining log transform by the
sum of independent cell log transforms with argument $p$. Each cell
variable has bare variance $O(r^2)$, so its subGaussian bound gives
\[
 \log\E e^{pX}=\tfrac12p^2\Var X+O_p(r^3),
\]
uniformly in the fine mesh. For example Taylor's integral remainder is
bounded by $C_p\E[|X|^3e^{p|X|}]=O_p(r^3)$ using the Gaussian tail.
There are $O(r^{-2})$ cells. Apply Equation~\eqref{st:upper-energy}, let
$n\to\infty$, then $r\downarrow0$, then $p\downarrow1$. This proves the
matching upper bound. Apply it to arbitrary linear combinations to obtain
the joint law; the uniform exponential bounds give convergence of moments.

Finally nearest-neighbor Poincar\'e gives
$\|u-\bar u\|_2^2\le C_Bn^2(u,A_{B_n}u)$. Consequently the bare variance
of a density error is bounded by its scaled $L^2$ norm squared. Approximate
$f$ by finitely many nonconstant Neumann cosines of the rectangle, and
use their smooth inverse-Laplacian shifts. For a Neumann shift,
$n^2A_{B_n}g$ converges in density $L^2$ to $-\Delta g$: interior errors
vanish by Taylor expansion; boundary coefficients are $O(n^{-2})$ because
the normal derivative vanishes and tangential steps cancel. At corners
both derivatives vanish. The $O(n)$ boundary-layer sites have vanishing
scaled area, so their bounded scaled errors disappear. Use the actual
discrete rectangle endpoints in this approximation and then let them
converge. Poincar\'e and bare domination complete the density approximation
and prove Equation~\eqref{st:rectangle-smear} and its joint version.
\end{proof}

A face probability can also be used in the zero-total smear conclusions.
Replacing it by a uniform strip of width $r$ costs at most $C_Jr+O(n^{-1})$:
send each row's mass normally over distance at most $rn$ with currents
$O(n^{-1})$. Apply Theorem~\ref{st:free-rectangle-limit} at fixed $r$ and
then let $r\downarrow0$. In particular, differences of face and volume
probabilities have bounded bare cost, and have variance tending to zero
when $a=0$.
\subsection{Neutral charges and the fractional mean}

For a free graph, $w_B(s)=\E e^{i(h,s)}$ is well defined independently of
the representative whenever $\sum_xs_x\in\Z$. Pin a vertex $o$ and apply
Poisson summation in the remaining $2\pi\Z$ coordinates. Extending a
residual profile by its opposite sum at $o$ identifies its quadratic form
with the free inverse-Laplacian form. The result is the exact identity
\begin{equation}\label{st:poisson-charge}
 w_B(s)=
 \frac{\displaystyle\sum_{q\in\Z^B:\,\sum q=\sum s}
       e^{-\frac k2(q-s,G_B(q-s))}}
      {\displaystyle\sum_{q\in\Z^B:\,\sum q=0}
       e^{-\frac k2(q,G_Bq)}}\in(0,1].
\end{equation}
There is no $2\pi$ in the dual energy, because the height spacing is $2\pi$.
Let $\nu_B$ be the centered neutral charge law in the denominator. Twice
differentiating Equation~\eqref{st:poisson-charge} at $s=0$ in directions
$A_Bg,A_Bg'$ gives
\begin{equation}\label{st:dual-covariance}
 \Cov_{\nu_B}((q,g),(q,g'))
 =k^{-1}(g,A_Bg')-k^{-2}\Cov_B((h,A_Bg),(h,A_Bg')).
\end{equation}
At an arbitrary integer-total profile, one differentiation gives
\begin{equation}\label{st:dual-mean}
 \partial_{A_Bg}\log w_B(s)
       =k(g,\E_s(q-s)).
\end{equation}
The dual Gaussian comparison implies that the Hessian of $\log w_B$ in
neutral directions is at least its centered Hessian, namely minus the
height covariance. Equation~\eqref{st:coefficient-transfer} applies here.

We will also use the following replica identity, with no limiting assumption.
For two independent centered height copies let $S=h_1+h_2$ and $D=h_1-h_2$.
Condition on their common parity class modulo $4\pi\Z$ after pinning.
Then $S,D$ are conditionally independent and identically distributed,
each with half the original precision. For linear tests $i,j$, put
$V_{ij}=\E[S_iS_j\mid\text{parity}]$. Expansion of the two independent
copies gives
\begin{equation}\label{st:replica-covariance}
 \E V_{ij}=2\Cov(h_i,h_j),\qquad
 \Cov(V_{ij},V_{rs})=2\kappa(h_i,h_j,h_r,h_s).
\end{equation}
Indeed $\E[V_{ij}V_{rs}]=\E[S_iS_jD_rD_s]$ by conditional independence;
the terms with an odd number of factors from either copy vanish, and the
remaining terms are $2\E[h_ih_jh_rh_s]+2C_{ij}C_{rs}
-2C_{ir}C_{js}-2C_{is}C_{jr}$. Subtract $4C_{ij}C_{rs}$.

\begin{theorem}[Uniform fractional mean]\label{st:fractional-mean}
Suppose $a(k)>0$, and let $u_m$ be the uniform probability on $B_m$.
Then $(h,u_m)$ modulo $2\pi$ converges to uniform measure on the circle,
independently of the Gaussian observations in
Theorem~\ref{st:free-rectangle-limit} on a square. More generally, for every
fixed $l\in\Z\setminus\{0\}$ and every sequence of zero-total profiles
$t_m$ of uniformly bounded bare inverse energy,
\begin{equation}\label{st:vanishing-phase}
 w_{B_m}(lu_m+t_m)\longrightarrow0.
\end{equation}
\end{theorem}

\begin{proof}
Write $g_m=(x_1-(m-1)/2)/m$, $X=(h,A_Bg_m)$, and
$Y=(h,\rho-u_m)$, where $\rho$ is the uniform right-face probability.
The flow bound makes $Y$ uniformly subGaussian. Reflection changes $X$
to $-X$ and exchanges right and left faces, so
\begin{equation}\label{st:face-covariance}
 \Cov(X,Y)\longrightarrow a/2.
\end{equation}
Here is a discrete verification of the normalization. The profile
$A_Bg_m$ is supported in fixed-width side layers. Except in boundedly many
corner sites, its right-layer row sums are constant, its left-layer row
sums are their negatives, and all other sums vanish, by pairing steps.
Transport these coefficients to their respective faces along bounded
paths; the cost is $O(m^{-1})$. Corner discrepancies of size $O(m^{-1})$
can be routed along $O(m)$ paths, at cost $O(m^{-1})$. Thus
$A_Bg_m-t_m(\rho_R-\rho_L)$ has vanishing flow cost. Pairing with $g_m$
and using $(g_m,A_Bg_m)\to1$ gives $t_m\to1$. This proves
Equation~\eqref{st:face-covariance}.

Fix $l\ne0$ and suppose along a subsequence
$w_B(lu_m)\to r>0$. Uniform subGaussian bounds on $X,Y$ allow a further
subsequence on which
\[
 F_m(t,u)=\E\exp\{il(h,u_m)+itX+iuY\}
\]
converges locally uniformly with all derivatives on $\C^2$.
This follows from the uniform bound
$|F_m(t,u)|\le\E e^{|\Im t||X|+|\Im u||Y|}$, Cauchy's formula, and a
diagonal Arzel\`a--Ascoli argument. The limit $F$ is real and nonnegative
on $\R^2$ by Equation~\eqref{st:poisson-charge}.

The Gaussian limit and moment convergence of $X$ imply
$\Var(V_{XX})=2\kappa(X,X,X,X)\to0$. The fourth moments of $Y$ are bounded,
so $\Var(V_{YY})$ is bounded. By Equation~\eqref{st:replica-covariance},
\[
 \Var(V_{XY})=2\kappa(X,X,Y,Y)
            =\Cov(V_{XX},V_{YY})\longrightarrow0.
\]
Multiply $V_{XX}$ or $V_{XY}$ by the conditional characteristic function
of $D$ with phase $l\bar D+tD_X+uD_Y$. Its modulus is at most one, so the
concentration just proved and conditional independence permit replacement
of $V_{ij}$ by $2C_{ij}$. Direct expansion in the two original copies gives
\[
 \E[S_iS_j e^{i(l\bar D+tD_X+uD_Y)}]
       =-2(F_m F_{m,ij}-F_{m,i}F_{m,j}).
\]
Since $\E e^{i(l\bar D+tD_X+uD_Y)}=F_m^2$, it follows wherever $F>0$ that
$(\log F)_{tt}=-a$ and $(\log F)_{tu}=-a/2$.
At $(0,0)$ reflection gives $F_t=0$. Analytic continuation yields
\begin{equation}\label{st:phase-transform}
 F(t,u)=F(0,u)\exp\{-at^2/2-atu/2\}.
\end{equation}
Equation~\eqref{st:coefficient-transfer} ensures $F(0,l)>0$, because the
bare cost of $lY$ is bounded.

Now join two horizontally adjacent $m$-squares to form a free rectangle,
and use its centered neutral charge law. For a charge $q$, let $Q$ be its
left total, the right total being $-Q$. Use the following trial quadratic
form: in each cell route the residual $q-Q_B\rho_B$ with its minimizing
internal flow, where $\rho_B$ is its seam-face probability; route $Q$
uniformly across the $m$ nearest-neighbor seam edges; add the penalty $Q^2$.
The resulting form is
\[
 k\sum_{B=L,R}(q_B-Q_B\rho_B,G_B(q_B-Q_B\rho_B))+D_mQ^2,
 \qquad D_m=1+k/(cm).
\]
It dominates the actual precision form $k(q,Gq)$ by
Equation~\eqref{st:flow}. Conditional on $Q$, the trial cells are independent
shifted neutral lattice Gaussians, whose affine variances dominate the
centered ones. By Equation~\eqref{st:dual-covariance} and saturation, their
sum tends to $2(k-a)/k^2$, exactly the actual variance of the global affine
test $x_1/m$ on the two-cell rectangle. Precision comparison gives the
opposite upper bound for the trial total variance. Thus the variance,
over $Q$, of its conditional affine mean tends to zero.

The sector weights, after common normalization, are
$e^{-D_mQ^2/2}w_L(Q\rho_L)w_R(-Q\rho_R)$.
Their normalizing sum is at most $\sum_Q e^{-Q^2/2}$ and the $Q=0$ term is
one. The $Q=l$ term stays positive by Equation~\eqref{st:coefficient-transfer}.
In that sector Equation~\eqref{st:dual-mean} and
Equation~\eqref{st:phase-transform} give
$k(g,\E[q-l\rho])\to-al/2$ in the left cell. Reflection and charge
reversal give the same residual pairing in the right cell. The two
seam-face contributions to the global affine mean differ by only
$O(m^{-1})$. Its $Q=l$ conditional mean therefore tends to $-al/k\ne0$,
whereas the $Q=0$ conditional mean is zero. This contradicts the vanishing
sector-mean variance. Thus $w_B(lu_m)\to0$.

If Equation~\eqref{st:vanishing-phase} failed, its positive lower bound
would transfer to $lu_m$ by Equation~\eqref{st:coefficient-transfer}, which
is impossible. Joint convergence with the Gaussian observations follows
by Fourier polynomials on the circle and joint characteristic functions;
tightness on the circle and the already proved Gaussian tightness make
this convergence of the full joint law. Exponential bounds extend the
joint assertion to moments of the Gaussian coordinates.
\end{proof}

\subsection{The charge gap and a uniform critical-scale ratio}

Let $\rho_{B,e}$ be uniform probability on face $e$ of a square, and set
$\bar\rho_B=\frac14\sum_e\rho_{B,e}$ and
$z_m=w_{B_m}(\bar\rho_{B_m})>0$.

\begin{lemma}[Face ratio with a common constant]\label{st:critical-face-ratio}
Fix $J,k$. Suppose every fixed nonzero integer-total bounded face profile
has characteristic coefficient tending to zero as the square side tends
to infinity. This hypothesis holds if $a(k)>0$. There are constants
$c_*=c_*(J,k)>0$ and an integer $G_0$, independent of $G$, such that for
\emph{every fixed integer} $G\ge G_0$,
\begin{equation}\label{st:face-ratio}
 \liminf_{m\to\infty}\frac{z_{Gm}}{z_m}
       \ge c_*G^{\,2-a(k)/(4\pi)}.
\end{equation}
The fine-size threshold implicit in this liminf may depend on $G$.
At $a(k)=8\pi$, the same positive lower bound $c_*$ therefore holds
for arbitrarily large fixed $G$.
\end{lemma}

\begin{proof}
Tile the $Gm$-square into $G^2$ cells $B_i$. The restriction of the large
square's face probability to $B_i$ has total $\xi_i$ and equals the sum of
its exterior-face probabilities with coefficient $1/(4G)$ per exterior
face. Thus $\xi$ is supported on boundary cells, $\sum_i\xi_i=1$, and
$\sum_i\xi_i^2\le C/G$.
For a neutral real profile $r$, let $Q_i=\sum_{B_i}r$ and choose a linear
map from $Q$ to an antisymmetric nearest-neighbor coarse flow $F$ satisfying
$\operatorname{div}F=Q$ and zero exterior flux. The quadratic form
\begin{equation}\label{st:trial-precision}
 \begin{split}
 \mathcal T_m(r)={}&k\sum_i
 \left(r_i-\sum_eF_{ie}\rho_{i,e},
 G_{B_i}\left(r_i-\sum_eF_{ie}\rho_{i,e}\right)\right)\\
 &+\frac{k}{cm}\sum_{\langle i,j\rangle}F_{ij}^2+\sum_iQ_i^2
 \end{split}
\end{equation}
is at least $k(r,G_{B_{Gm}}r)$: use the internal minimizing flows and
send each coarse flux uniformly across the $m$ nearest-neighbor seam
edges. The last term is an extra positive penalty. The internal and seam
edges are disjoint, so their costs add. The form is positive definite on
the neutral space because zero cost forces $Q=0$ and then $r=0$.

In the numerator of Equation~\eqref{st:poisson-charge} take
$r=q-\bar\rho_{B_{Gm}}$; in the denominator take $r=q$. Precision
monotonicity of a shifted-to-centered ratio shows that replacement by
$\mathcal T_m$ gives a lower bound for $z_{Gm}$.
Sum this new partition function by integer cell totals $n_i=\sum_{B_i}q$.
Then $Q=n-\xi$ in the numerator and $Q=n$ in the denominator.
Divide both sums by the product of the cells' centered neutral partition
functions. Their internal factors become $\prod_iw_{B_i}(s_i)$, where
\[
 s_i=\sum_eH_{ie}\rho_{i,e},\qquad
 H_{ie}=\begin{cases}
 F_{ie},&\text{internal face},\\
 1/(4G),&\text{exterior face in the numerator},\\
 0,&\text{exterior face in the denominator}.
 \end{cases}
\]
In the denominator exterior $H$ is zero. In either case $\sum_eH_{ie}=n_i$.
All these factors lie in $(0,1]$. The penalty
$\exp(-\frac12\sum_iQ_i^2)$ is summable over the finite-rank integer sector
lattice, uniformly in $m$. Thus dominated convergence is valid at every
fixed $G$. Each nonzero denominator sector contains a nonzero integer-total
face profile and tends to zero by hypothesis. The zero sector equals one,
so the normalized denominator tends to one.

Keep numerator sectors $n=e_{i_0}$ with $i_0$ at distance at least $G/4$
from the boundary. Write $Y_i=s_i-n_i\bar\rho_{B_i}$. The logarithmic Hessian
bound for $w$ gives
\begin{equation}\label{st:face-product}
 \prod_iw_{B_i}(s_i)
 \ge z_m\exp\left\{-\tfrac12\sum_i\Var_{B_i}(h,Y_i)\right\}.
\end{equation}
For the linear term at zero, inversion symmetry suffices. On the total-one
hyperplane put
$\Psi_B(H)=\log w_B(\sum_eH_e\rho_{B,e})$, where $\sum_eH_e=1$.
Equation~\eqref{st:poisson-charge} makes this function smooth there, and
square symmetry makes its differential at $H_0=(1/4,1/4,1/4,1/4)$
invariant under the square group. Every tangent vector with zero coordinate
sum has group average zero, so this differential vanishes on it. The face
coefficients of each $Y_i$ have zero sum, and $n=e_{i_0}$ uses only the
base profiles zero and $\bar\rho_{B_i}$.

Put $D_{ij}=(H_{i,+j}-H_{i,-j})/2$ and
$S_{ij}=H_{i,+j}+H_{i,-j}$. Reflection in each coordinate splits $Y_i$ into
two odd face differences and an even remainder. These three pieces have
zero mutual covariance by symmetry. The odd differences have limiting
variance $a$ and zero cross covariance, as in
Equation~\eqref{st:face-covariance}; the even remainder has a bounded bare
variance times $\sum_jS_{ij}^2$, by the face-to-strip flow estimate.
Consequently, for fixed $G$ and fixed face coefficients,
\begin{equation}\label{st:face-energy-bound}
 \limsup_m\sum_i\Var(h,Y_i)
 \le a\sum_{i,j}D_{ij}^2+C_{J,k}\sum_{i,j}S_{ij}^2.
\end{equation}

We construct one coarse flow for each retained sector with bounds uniform
in $G,i_0$. In the continuum square of side $G$, let $x_0$ be the center
of cell $i_0$. There is a vector field with divergence $\delta_{x_0}$ and
constant outer normal flux $1/(4G)$, of the form
\begin{equation}\label{st:radial-flow}
 V(x)=\frac{x-x_0}{2\pi|x-x_0|^2}+E(x),\qquad
 |E(x)|\le C/G,\quad |\nabla E(x)|\le C/G^2.
\end{equation}
For completeness, subtract the expanding field
$(x-x_{\rm center})/(2G^2)$, which has divergence $G^{-2}$ and the desired
boundary flux. Solve the resulting zero-Neumann problem by even reflection
to the torus of side $2G$. After rescaling to a fixed torus, subtract
cutoff logarithms around its finitely many reflected poles. Their mutual
separations and distances from the original boundary are bounded below,
uniformly for $x_0/G$ in the central square. The remainder solves a smooth
mean-zero Poisson equation with uniformly bounded derivatives. Its Fourier
series has coefficients $\widehat f(p)/|p|^2$; repeated integration by
parts gives summability after any fixed number of differentiations.
This proves the smooth bounds in Equation~\eqref{st:radial-flow}, including
up to the reflected boundary. Adding back the expanding field proves the
assertion.

Integrate $V$ over the faces of each continuum unit cell, with outward
sign, to define $H$. It has divergence $e_{i_0}$ and the prescribed exterior
flux. The central faces have bounded fluxes. Off the central cell,
\[
 |S_{ij}|\le C(1+|i-i_0|)^{-2}+C/G^2,
 \qquad D_{ij}=V_j(i)+O((1+|i-i_0|)^{-2}+G^{-2}).
\]
Hence
\begin{equation}\label{st:radial-cost}
 \sum_{i,j}S_{ij}^2\le C,\qquad
 \sum_{i,j}D_{ij}^2\le\frac1{2\pi}\log G+C.
\end{equation}
The coefficient follows by integrating the squared radial field:
$\int_1^{CG}(4\pi^2r^2)^{-1}2\pi r\,dr=(2\pi)^{-1}\log G+O(1)$.
The comparison of the sum and integral has summable derivative error.
The cross term with $E$ is $O(G^{-1}\sum_{r\le CG}1)=O(1)$, and
$\sum|E|^2=O(1)$. All constants are uniform for the retained cells.
The internal part $F$ has divergence $e_{i_0}-\xi$. These vectors, for
central $i_0$, are linearly independent because their central coordinates
are distinct and $\xi$ is supported at the boundary. Prescribe the linear
flow map on them by this construction and extend it to a basis of the
neutral coarse space, using any path flow on the remaining basis vectors.
Thus a single quadratic form~\eqref{st:trial-precision} serves all retained
sectors; separate nonlinear choices of precision have not been made.

For these sectors the penalty $\|e_{i_0}-\xi\|_2^2$ is bounded by a common
constant. The crossing cost tends to zero as $m\to\infty$ for each fixed
$G$. Equations~\eqref{st:face-product}--\eqref{st:radial-cost} therefore give
at least $c z_mG^{-a/(4\pi)}$ from each sector in the liminf of the ratio.
There are at least $c'G^2$ retained sectors. Division by the denominator,
which tends to one, proves Equation~\eqref{st:face-ratio}, with one positive
constant for all $G\ge G_0$.
\end{proof}

\begin{corollary}[Exact forbidden range]\label{st:charge-gap}
For every admissible $J$ and $k>0$,
\begin{equation}\label{st:gap}
 a(k)=0\quad\hbox{or}\quad 8\pi\le a(k)\le k.
\end{equation}
\end{corollary}

\begin{proof}
If $0<a<8\pi$, Theorem~\ref{st:fractional-mean} gives the hypothesis of
Lemma~\ref{st:critical-face-ratio}. Choose a fixed $G$ so large that its
right side exceeds two. The definition of liminf then gives
$z_{Gm}\ge\frac32z_m$ for all sufficiently large $m$. Iteration along
$m,Gm,G^2m,\ldots$ contradicts $0<z_m\le1$. The upper bound $a\le k$ was
proved in Theorem~\ref{st:free-rectangle-limit}.
\end{proof}

\subsection{Localization when the structural coefficient vanishes}

\begin{proposition}[Zero coefficient implies localization]\label{st:zero-localization}
If $a(k)=0$, there is a finite constant $C_{J,k}$ such that
\begin{equation}\label{st:localized-criterion}
 \sup_{R\ge1}\E_{\mu^0_{J,v^2k,R}}h_0^2\le C_{J,k}.
\end{equation}
Moreover, there is a fixed integer $m\ge1$ for which centered height
laws on every sufficiently large square torus of side divisible by $m$
satisfy $\sup_{x,y}\E(h_x-h_y)^2\le C_{J,k,m}$.
\end{proposition}

\begin{proof}
\emph{Good phases at arbitrarily large scales.}
Some nonzero mean Fourier coefficient stays positive along a subsequence.
Otherwise $w_{B_m}(ju_m)\to0$ for all $j\ne0$. Every fixed face profile of
integer total $j$ differs from $ju_m$ by a test of variance tending to
zero, by the face extension of Theorem~\ref{st:free-rectangle-limit} at
$a=0$. The inequality
$|\E e^{iU}-\E e^{iV}|\le\sqrt{\E(U-V)^2}$ would then give the hypothesis
of Lemma~\ref{st:critical-face-ratio} with $a=0$, whose growth contradiction
is impossible. Hence along a subsequence, extracted diagonally over $j$,
\[
 w_{B_m}(ju_m)\longrightarrow r_j=r_{-j}\ge0,
 \qquad r_0=1,\qquad r_l>0\quad\text{for some }l\ge1.
\]

Use Equation~\eqref{st:trial-precision} on a $Gm$-square, now with shift
$l u_{Gm}$ instead of the face probability. Cell totals are
$Q_i=n_i-l/G^2$ in numerator sectors, with $\sum_i n_i=l$, and $Q_i=n_i$
in denominator sectors. For each fixed $G$ and sector, the crossing cost
vanishes as $m\to\infty$. The internal shifted profile differs from
$n_i u_m$ by a fixed linear combination of face and volume probabilities
of total zero, so its coefficient tends to $r_{n_i}$. The quadratic
penalty gives summable domination of the sector sums, exactly as in
Lemma~\ref{st:critical-face-ratio}. Therefore, with $b(j)=r_je^{-j^2/2}$,
\begin{equation}\label{st:phase-amplification}
 \liminf_m w_{B_{Gm}}(lu_{Gm})
 \ge e^{l^2/(2G^2)}\frac{b^{*G^2}(l)}{b^{*G^2}(0)}.
\end{equation}
The ratio on the right tends to one as $G\to\infty$. Here is a direct
local limit argument. Normalize $b$ to a probability $p$. It is even, has
exponential tails, positive mass at zero, and positive variance since
$p(l)>0$. Let $d\Z$ be its support subgroup; then $l\in d\Z$. Its
characteristic function $\psi$ has modulus less than one away from the
finitely many points $2\pi j/d$ in a period: equality in the triangle
inequality, together with $p(0)>0$, requires all support phases to equal
one. Near each such point,
$\psi(t)=1-\frac12\sigma^2(t-t_0)^2+O(|t-t_0|^4)$.
Fourier inversion of $\psi(t)^N e^{-ilt}$, splitting into fixed small
neighborhoods of these points and their complement, gives
$p^{*N}(l)=d(2\pi N\sigma^2)^{-1/2}(1+o(1))$, and the same formula at zero.
The complement is exponentially small, and the neighborhood calculation
uses $u=\sqrt N(t-t_0)$ with a Gaussian dominating bound. This proves the
ratio claim. Choose $G$ first and then arbitrarily large members $m$ of
the subsequence in Equation~\eqref{st:phase-amplification}. We obtain
arbitrarily large scales at which $w(lu)$ is arbitrarily close to one.

\emph{A joint phase bound.}
Fix $0<d_0<\pi/10$. Tile a torus by $m$-squares, with averages $\bar h_i$.
For any $\epsilon>0$ we can choose one of the preceding good scales so
that, for every collection $I$ of distinct cells,
\begin{equation}\label{st:phase-bad}
 \Prob\{\dist(l\bar h_i,2\pi\Z)\ge d_0\text{ for all }i\in I\}
 \le\epsilon^{|I|},
\end{equation}
uniformly in the number of torus cells. We provide the comparison that
makes this a joint, rather than single-cell, assertion.
Add independent $N(0,s^2)$ noises to the phases and wrap them modulo
$2\pi$. Let $p_T(y)$ be the resulting joint density. Expanding the wrapped
normal introduces integers $j_i$ and the quadratic costs
$\sum_i(y_i+2\pi j_i-l\bar h_i)^2/(2s^2)$.
Together with the height energy these are a lattice Gaussian on the
height and $j$ coordinates, shifted only in the $j$ coordinates. Add a
small centered height mass to make it positive definite.
Increasing the height precision from the cut free cells to the torus
therefore decreases $p_T(y)/p_T(0)$ by
Lemma~\ref{st:finite-gaussian}. Let the mass decrease to zero using the
translation-invariant clause of that lemma. This yields
\[
 \frac{p_T(y)}{p_T(0)}\le\prod_{i\in I}\frac{p_i(y_i)}{p_i(0)}.
\]
The wrapped-noise supremum bounds $p_T(0)\le(C/s)^{|I|}$.
Fourier positivity in Equation~\eqref{st:poisson-charge} bounds
$p_i(0)\ge(2\pi)^{-1}$. At a good scale the free-cell phase is within any
fixed small distance of zero with probability arbitrarily close to one:
use $\E\cos(l\bar h)=w(lu)$ and
$1-\cos x\ge1-\cos\delta$ when $\dist(x,2\pi\Z)\ge\delta$.
Thus, for the noisy free-cell phase,
\[
 \Prob\{\dist(l\bar h+Z,2\pi\Z)\ge d_0/2\}
       \le C e^{-c d_0^2/s^2}+\delta_m,
 \qquad \delta_m\longrightarrow0
\]
along increasingly good scales. Integrate the density-ratio bound over
the product of these bad sets. Conditional on original phases which are
all bad at threshold $d_0$, the independent noises retain badness at
threshold $d_0/2$ with probability at least
$\Prob(|Z|\le d_0/2)^{|I|}$. Hence the left side of
Equation~\eqref{st:phase-bad} is at most
\[
 \left[
 \frac{C}{s\Prob(|Z|\le d_0/2)}
       \bigl(e^{-c d_0^2/s^2}+\delta_m\bigr)
 \right]^{|I|}.
\]
Choose $s$ small first, then a sufficiently good large scale. This proves
Equation~\eqref{st:phase-bad} with the stated uniformity.

\emph{Neighbor differences and bad animals.}
Declare a coarse cell bad if its phase is bad, or if its average differs
from any axis-neighbor average by more than $d_0/l$.
For any list of $q$ pairs of neighboring cells with mutually disjoint
end cells, precision comparison after cutting into pair rectangles gives,
for each choice of signs $\eta_j\in\{-1,1\}$,
\[
 \E\exp\left\{T\sum_{j=1}^q\eta_j
                   (\bar h_{i_j}-\bar h_{i'_j})\right\}
 \le e^{q\delta_m(T)},\qquad \delta_m(T)\longrightarrow0.
\]
The free pair difference is a zero-total bounded density test, so the
limit follows from Equation~\eqref{st:rectangle-smear} at $a=0$; the bound
is uniform in the pair's location and orientation. Chernoff and the
$2^q$ sign choices give a bound
$[2e^{-Td_0/l+\delta_m(T)}]^q$ for simultaneous failures. Choose $T$ large
first and $m$ large afterward. This exponential base can thus be made as
small as desired, simultaneously with Equation~\eqref{st:phase-bad}.

For a prescribed sup-norm-connected set of $j$ coarse sites, assign each
bad site one of its five failure types (phase or one of four directed
neighbors). One type occurs at least $j/5$ times. For an edge type, a
greedy extraction gives at least $j/C_0$ mutually disjoint pairs, since
each pair meets only a bounded number of pairs of the same type. For the
phase type use Equation~\eqref{st:phase-bad} directly. Union over the at
most $5^j$ assignments. Taking the preceding bases sufficiently small
proves, with $\eta>0$ arbitrarily small,
\begin{equation}\label{st:animal}
 \Prob\{\text{every site of a prescribed sup-connected size-$j$ set is bad}\}
 \le\eta^j.
\end{equation}
The number of such sets containing a specified root is at most $C_1^j$:
a depth-first traversal of a spanning tree has at most $2j$ steps, each
with at most eight choices. We fix the scale so that $C_1\eta$ is smaller
than a sufficiently small absolute constant. Good axis-neighbors have
the same nearest integer label to $l\bar h/(2\pi)$, because the difference
between their labels times $2\pi$ is at most $3d_0<2\pi$.

\emph{Uniform second moments from coarse connections.}
Let the coarse torus side be $M$. For any two prescribed coarse cells and
$1\le t\le M/200$, their radius-$Ct$ neighborhoods are joined by a good
nearest-neighbor path except on an event of probability
\begin{equation}\label{st:connection-tail}
 C e^{-c_0t},
\end{equation}
where $C,c_0>0$ do not depend on $M$ or on the two cells.
Here are the geometric and counting details. A missing good crossing of
a coarse rectangle requires a sup-connected bad crossing in the
perpendicular direction: take the boundary of the set reachable from the
starting side by good nearest-neighbor paths and trace its separating
boundary. Its length is at least the shorter rectangle side. For bounded
aspect ratio, Equation~\eqref{st:animal} and the preceding rooted count
bound its probability by $Cs^2e^{-cs}$ when the width is comparable to $s$.
Four overlapping strip crossings produce a good circuit in a square
annulus of radii $s,2s$. Request such circuits around each prescribed cell
for $s=t,2t,4t,\ldots$ up to one scale $s_*$ comparable to $M/100$.
Join consecutive circuits by a crossing of a bounded-aspect rectangle
whose ends are inside the smaller circuit and outside the larger.
Connect the two largest circuits by a bounded number of circuits and
crossing rectangles along a chain of centers spaced by at most $s_*/4$.
All these rectangles have lifts to planar rectangles of bounded aspect,
so the same separation argument applies even on the torus. A union bound
sums at most $C\sum_{r\ge0}(2^rt)^2e^{-c2^rt}$, plus a bounded number of
$M$-scale errors, and proves Equation~\eqref{st:connection-tail} after
adjusting constants. For $t$ comparable to $M/200$ one uses the same
construction with a single macroscopic scale.

Write $A_i=\bar h_i$. On successful connection there are good end cells
$i',j'$ within distance $Ct$ of the prescribed cells $i,j$, and
$|A_{i'}-A_{j'}|\le2d_0/l$ by label agreement. Let
\[
 B_t=2d_0/l+
 \max_{\dist(i,i')\le Ct}|A_i-A_{i'}|
 +\max_{\dist(j,j')\le Ct}|A_j-A_{j'}|.
\]
Nearest-neighbor trial paths and the bare Gaussian estimate imply
$\E B_t^4\le C_m(1+t)^p$ for a fixed integer $p$ independent of $M$.
Explicitly, connect any two vertices of cells at coarse distance at most
$Ct$ by a path of length at most $C m(t+1)$. Its flow cost is at most
$C_Jm(t+1)$; averaging such paths bounds the average-difference variance
by that number. Its fourth moment is at most a constant times the square
of that variance, by the subGaussian tail. There are $O((1+t)^2)$ possible
end cells for each maximum, which proves a polynomial bound, for example
with $p=4$. The constants may depend on the now fixed $m$.

Partition into the first successful integer value of $t$. For $t\ge2$,
that event is contained in failure at $t-1$, whose probability is at most
$Ce^{-c_0(t-1)}$. Cauchy--Schwarz therefore gives
\[
 \E(A_i-A_j)^2\ind_{\{\text{first success at }t\}}
 \le C_m(1+t)^{p/2}e^{-c_0(t-1)/2}.
\]
The $t=1$ term is bounded by the same moment bound without its exponential
factor. Sum until the final macroscopic scale. On remaining failure,
use a torus path of length $O(mM)$ to bound the fourth moment of
$A_i-A_j$ by $C_mM^2$ and multiply by the square root of its exponentially
small failure probability. This term is uniformly bounded. Thus
$\sup_{M,i,j}\E(A_i-A_j)^2<\infty$ for large $M$. Within a fixed cell,
the point-minus-average fourth moment is uniformly bounded by paths of
length $O(m)$, so the same uniform second-moment bound holds for every
pair of prescribed vertices. Finitely many smaller coarse tori can be
absorbed into the constant.

Finally embed any desired $Q_R$ and its $J$-interaction collar in a much
larger torus with side divisible by $m$. Choose a vertex $y$ outside that
collar and impose the centered equalities $h_z=h_y$ at every vertex outside
$Q_R$. The conditional relative law on $Q_R$ is exactly the zero-exterior
law at temperature $v^2k$; no interaction wraps across the embedding.
Centered equality comparison in Lemma~\ref{st:finite-gaussian} bounds
$\Var(h_0-h_y)$ in this conditional law by its unconditioned torus variance.
The latter has the bound just proved, independent of $R$. Symmetry centers
the conditional variable and proves Equation~\eqref{st:localized-criterion}.
\end{proof}
